\section{Comparison of Existing Platforms}\label{sec:comparison-of-existing-platforms}
This section compares two widely used web platforms, WordPress and Squarespace.

\subsection{WordPress}\label{subsec:wordpress}
WordPress~\parencite{wordpress} is a widely adopted open-source content management system written in PHP\@.
It is flexible through its plugin and theming ecosystem.
The platform is open-source, which enables the developers to modify the core behavior if needed.

The primary advantages of WordPress include:

\begin{itemize}
    \item High flexibility due to being open-source.
    \item Large ecosystem of plugins and themes.
    \item Strong community support and documentation.
\end{itemize}

Excessive plugin usage may introduce inconsistent behavior or even produce security problems.
However, this flexibility comes at the cost of architectural complexity, which is partially addressed by code comments
and in-depth documentation.

From a developer perspective, WordPress is powerful but suffers from complex architecture accumulated over the years
of development~\parencite{wordpress_docs_rest}.

\subsection{Squarespace}\label{subsec:squarespace}
Squarespace~\parencite{squarespace} is a proprietary, fully hosted website builder focused on ease of use and visual design.
Even though Squarespace is closed-source system, it provides viable plugin system.

It may seem like closed-systems are more secure over the open-source projects.
However, this is not necessarily the case.
When security flaw is found in open-source system it is either fixed before it has a chance to be exploited
or the issue is fixed before it can negatively impact large number of users.
For example, the backdoor in xz Utils was identified through external analysis from third-party developer,
whose code depends on the library~\parencite{cve_2024_3094}.
This is the reason why the thesis proposed solution is open-source.

The main advantages of Squarespace include:

\begin{itemize}
    \item Fast setup and deployment.
    \item Integrated hosting and maintenance.
    \item Consistent user experience and stability.
    \item Minimal technical knowledge required.
\end{itemize}

The primary limitation of Squarespace is restricted extensibility of the core systems.
Developers cannot access the underlying backend architecture to host the websites on their own servers or to modify
core behavior.

Squarespace design philosophy prioritizes end users over developers,
which is reflected by the closed-source nature.
This makes it well-suited solution for less tech-savvy users, but lacks in flexibility for developers.

\subsection{Platforms Comparison Summary}\label{subsec:platforms-comparison-summary}
WordPress emphasizes openness and extensibility but suffers from architectural complexity.
Squarespace prioritizes simplicity but lacks customization needed for developers.
From a usability perspective, the Squarespace website editor is easier to use than the one WordPress provides.

Neither platform fully satisfies the requirements of a system that aims to provide:

\begin{itemize}
    \item Strong foundational architecture and modern code base.
    \item Full control over backend behavior.
    \item Extensibility without dependency-heavy ecosystems.
\end{itemize}

The thesis solution positions itself between these two approaches:
maintaining developer control and being open-source project similar to WordPress
while providing easy to use User Interface like the one Squarespace defines.
